% =====================================================================
%  GENERAL INTRODUCTION
% =====================================================================

\chapter*{General Introduction}
\addcontentsline{toc}{chapter}{General Introduction}
\markboth{General Introduction}{General Introduction}

% ---------- 1. THE PROBLEM ----------

Sooner or later, every large institution has to write down who is allowed to
decide what. In a development bank that question has teeth. An approval signed
by the wrong person is a governance failure, and auditors find those. The
African Development Bank Group keeps its answer in a document called the
Delegation of Authority Matrix, the DAM, which goes through the Bank's
activities one by one and records who starts each task, who has to check it, who
reviews it, who approves it, and who simply has to be told it happened.

The document is accurate. Getting at what it says is the problem.

The full DAM runs past two hundred and fifty pages. Nearly all of it sits inside
wide tables, and because those tables carry more columns than a page can hold,
the role names at the top are printed sideways. Authority is written in a
shorthand of single letters, and a small digit beside a letter can change what
it means. So a staff member who wants to know whether they can sign off on a
mission programme has to find the right row, in the right table, in the right
chapter, follow a rotated header down to work out which column belongs to which
role, and then decode the character where the two meet. People do manage this.
They manage it slowly, and not always correctly.

% ---------- 2. THE CLIENT ORGANISATION ----------

The African Development Bank Group was founded in 1964 and is headquartered in
Abidjan, C\^ote d'Ivoire. It has 81 member countries, 54 of them African and 27
from outside the continent, and its purpose is to support economic development
and social progress across its regional member countries. The Group also
includes the African Development Fund, which lends on concessional terms to the
countries least able to borrow commercially, and the Nigeria Trust Fund.

A bank that size cannot be casual about who decides things. It moves public and
multilateral money across dozens of countries, and every operational act, from
approving a mission to committing funds against a country programme, has to
trace back to somebody entitled to perform it. The DAM is how the Bank writes
that down, and it is where this project starts.

The work was carried out within the Master's programme in Data Sciences at
Polytech Intl. The African Development Bank Group acted as the client, supplying
both the source document and the requirements the finished system had to meet.

% ---------- 3. WHY IT MATTERS TO THE BANK ----------

None of the Bank's priorities get delivered unless decisions are made correctly,
by people entitled to make them. Take RDGN, the regional office covering
Tunisia, Algeria, Morocco, Mauritania and Libya from a hub in Tunis. Staff there
spend their days on exactly the activities the Matrix governs: organising
missions, dealing with co-financiers, committing funds against country
programmes. Each one needs an answer to the same four questions. Who starts it,
who checks it, who approves it, who needs to know. Delegation of authority is
not paperwork sitting alongside the real work. It is how the real work gets
done, and anything that makes it harder to consult slows the Bank down.

% ---------- 4. WHAT THE THESIS DELIVERS ----------

The obvious thing to try is a chatbot. It does not work here, and the reason is
worth being precise about. Language models write fluently whether or not they
are right, and fluent text about who controls financial approval is dangerous
when it is wrong. A model that invents a plausible-sounding role, or quietly
leaves out one of the three people who should have been informed, hands you
something that reads exactly like a correct answer. For a compliance document,
that is worse than no answer.

So the system described here is built around that problem rather than hoping to
avoid it. The DAM is parsed into a structured knowledge base. That knowledge
base becomes a graph of activities and the roles responsible for them, and
questions are answered by pulling real records out of that graph. There is a
language model involved, but it never decides what the facts are. It is handed
facts that have already been retrieved and asked to put them into a sentence,
and before anyone sees that sentence it is checked, word for word, against the
facts it came from. If a role name has gone missing or changed, the sentence is
thrown away and a plain one is shown instead.

The Bank's professional supervisors saw the system and their first reactions
were positive. It has not been through a full operational evaluation, so this
thesis makes no claim about measured time savings. What it can say is what the
system takes away: the hunt through two hundred and fifty pages, the rotated
header, the decoded letter. A question now comes back in seconds, with the
activity it came from attached. Whether that adds up to real efficiency gains
across the regional office is a question for a deployment study, not for this
report.

% ---------- 5. ORGANISATION OF THE THESIS ----------

The report runs to five chapters.

\textbf{Chapter 1} sets the scene: the host institution, the client, how the
Matrix is consulted today, where that breaks down, what existing tools offer,
and the requirements and architecture that follow from all of it.

\textbf{Chapter 2} looks hard at the source document. There is no tidy dataset
here, only a PDF whose structure lives in its layout, so this chapter works
through how that layout carries meaning and where automated extraction goes
wrong on it.

\textbf{Chapter 3} is the pipeline that turns the document into data, stage by
stage, along with the checks applied at each one and the knowledge graph built
on the output.

\textbf{Chapter 4} covers the agent: how a question finds the right activity,
how an answer is assembled, what the language model is allowed to do, and the
check that keeps it honest.

\textbf{Chapter 5} presents what was actually delivered, including the features
added in the project's second phase, the analytics and forecasting layer, the
tests behind it, and how it is deployed. The general conclusion then takes stock
of what worked, what did not, and what should be tackled next.
